\documentclass[11pt,a4paper]{article}
\usepackage[T1]{fontenc}\usepackage[utf8]{inputenc}\usepackage{mathpazo}
\usepackage[margin=2.2cm]{geometry}
\usepackage{booktabs,tabularx,array,multirow,graphicx,caption,float,enumitem,amsmath,amssymb,xcolor,fancyhdr}
\usepackage[colorlinks=true,linkcolor=blue!50!black,urlcolor=blue!50!black]{hyperref}
\newcolumntype{Y}{>{\raggedright\arraybackslash}X}
\renewcommand{\arraystretch}{1.18}\setlength{\parindent}{0pt}\setlength{\parskip}{4pt}
\setlength{\headheight}{14pt}\setlength{\emergencystretch}{3em}
\pagestyle{fancy}\fancyhf{}\lhead{\itshape AI251P Applied Computer Vision}\rhead{\itshape Phase 2 -- Project 5}\cfoot{\thepage}
\newcommand{\fig}[3]{\begin{figure}[H]\centering\includegraphics[width=#2\linewidth]{figures/#1}\caption{#3}\end{figure}}
\begin{document}
\begin{center}
{\large AI251P Applied Computer Vision}\\[4pt]
{\LARGE\bfseries Capstone Project -- Phase 2 Report}\\[6pt]
{\Large Automated Triage of Camera-Trap Imagery}\\[4pt]
Project 5\;\textperiodcentered\; Four members\;\textperiodcentered\; M4, M5, M8 (45 points)\\[3pt]
Submission deadline: 22 November 2026
\end{center}

\section{Team and system}
\par\noindent\begin{tabularx}{\linewidth}{@{}l Y@{}}\toprule
Field & Project information\\\midrule
Topic & Project 5: ecology and conservation; camera-trap event triage\\
Repository & \url{https://github.com/salehdev3/Capstone_5/tree/phase-2}; report-input evidence baseline \texttt{8ffeb2d}\\
W\&B project & \href{https://wandb.ai/saleh-alanazi-ministry-of-interior/camera-trap-phase2}{camera-trap-phase2}\\
Evidence backup & \href{https://wandb.ai/saleh-alanazi-ministry-of-interior/camera-trap-phase2/runs/mns1u8vt}{Run mns1u8vt}; \texttt{phase2-recovery-evidence:v0}\\
Service / model card & Local FastAPI batch endpoint; release/model-card details in Section~4\\\bottomrule
\end{tabularx}\par
\subsection*{Team members and ownership}
\par\noindent\begin{tabularx}{\linewidth}{@{}l l l Y@{}}\toprule
Name & Student ID & Workstream & Phase 2 contribution\\\midrule
Saud Abdullah Alfaiz & 225462 & Data/classification & Frozen data, 100-image box review; classifier recovery workflow\\
Saleh Mohammed Alanazi & 225486 & Detection & Animal detector, operating point and detector error review\\
Yousef Nafi Alanazi & 225457 & Segmentation & 230-mask review and segmentation experiments\\
Mohannad Mohammed Alqahtani & 225461 & Integration & HTTP service, timing, acceptance audit; v2 failure-review confirmation\\\bottomrule
\end{tabularx}\par
Each member owns one workstream. The five-member open-world extension is not applicable.

\subsection*{System overview}
Wildlife reviewers submit images with IDs, event IDs and timestamps. Animal boxes support species classification; masks provide a relative visible-body-area proxy. Outputs are an event table and a review gallery. Event count is the largest simultaneous detector count among retained frames, not the number of unique individuals.
\fig{pipeline.png}{1}{Image-folder pipeline. Detection/mask inputs: 640 pixels; species crops: $224\times224$. M4 counts use confidence .045/NMS .30. The mask helper repeats M4 inference; its cost is included in timing. Mask acceptance is specified in Section~3.3.}
\par\noindent\begin{tabularx}{\linewidth}{@{}l l Y Y@{}}\toprule
Component & Model & Input / output & Function\\\midrule
Classifier & Frozen ResNet-50 & Crop $\to$ 14 species probabilities & Species and review flag\\
Detector & YOLO11n & Image $\to$ animal boxes & Triage and counts\\
Segmentation & YOLO11n-seg + SAM & Predicted boxes/points $\to$ masks & Relative area and review\\\bottomrule
\end{tabularx}\par

\subsection*{Phase 1 continuity and evaluation scope}
The submitted Phase 1 report is unchanged. It selected frozen ResNet-50 on validation macro-F1 and reported species test accuracy .742 and macro-F1 .482. Its executable classifier checkpoint was absent from the available archives. A \emph{new} frozen ResNet-50 was trained for Phase 2 with the same 14-class vocabulary, split and core recipe. New test accuracy is .742126 and macro-F1 .482517. Similar metrics do not imply identical weights; the checkpoint hash and evidence are separately versioned.

The frozen manifest contains 6,037 frames across camera/event-disjoint train, selection-validation, calibration and test partitions. Camera, event and decoded-pixel groups are checked across splits. This is a project-specific CCT20 subset, not the official benchmark. Detection excludes incomplete positive-box frames. Segmentation uses 150 train, 40 selection-validation and 40 test images. Recovery v2 followed inspection of v1 test errors: its test result is \textbf{development evidence, not an independent generalization estimate}. The classifier test is historical, already reported in Phase 1. No new external unseen evaluation is claimed.

\subsection*{Acceptance criteria}
The four Phase 1 criteria remain unchanged. Three application criteria were added in the Phase 2 protocol before detector training; these additions are explicit, not retrospectively relaxed.
\par\noindent\begin{tabularx}{\linewidth}{@{}Y l l Y@{}}\toprule
Criterion & Metric & Target & Conditions\\\midrule
Preserve animals (P1) & Frame recall & $\geq .99$ & Retained test frames; selected detector point\\
Remove empties (P1) & Empty removal & $\geq .20$ & 53 retained empty test frames\\
Species (P1) & Macro-F1 & $\geq .65$ & 14 classes, 508 historical crops\\
Calibration (P1) & ECE & $\leq .10$ & 15 bins; temperature from calibration split\\
Mask overlap (added P2) & Mean IoU & $\geq .70$ & 40 reviewed development-test images\\
Counts (added P2) & Event MAE & $\leq 1$ & Complete supervision on every retained frame\\
Latency (added P2) & p95 & $\leq 2$ s & CPU, batch 1, fixed validation cases\\\bottomrule
\end{tabularx}\par

\subsection*{Follow-up on the Phase 1 data-improvement plan}
\par\noindent\begin{tabularx}{\linewidth}{@{}Y l Y@{}}\toprule
Planned action & Status & Evidence / measured effect\\\midrule
Preserve grouped split and audit labels & Done & Frozen manifest checks; Saud verified 100 training-frame reviews\\
Expand rare-class / empty-frame coverage & Not done & No new independent collection; species macro-F1 .483 and animal recall below target\\
Improve mask supervision & Partly & 230 reviewed references; six background polygons removed in a training-only copy. Continuation plus cleanup raises raw validation IoU .488 to .633; effects are not isolated\\
Analyse lighting and small objects & Partly & Grayscale groups and five small objects evaluated; verified IR/night labels unavailable\\\bottomrule
\end{tabularx}\par

\section{Object detection (M4)}
\textbf{Task.} Detect one \texttt{animal} category from human CCT20/team boxes. Species recognition belongs to the separate classifier; detector AP is not species accuracy.

\subsection{Annotation conversion and verification}
Normalized corners become YOLO class/centre/size labels:
\[(x_1,y_1,x_2,y_2)\mapsto\left(0,\frac{x_1+x_2}{2},\frac{y_1+y_2}{2},x_2-x_1,y_2-y_1\right).\]
Finite coordinates, positive dimensions and human annotation source are asserted. Negative frames have empty labels. Coordinates are preserved rather than silently clipped.
\par\noindent\begin{tabularx}{\linewidth}{@{}Y r Y@{}}\toprule
Item & Count & Notes\\\midrule
Retained images / boxes & 5,242 / 4,934 & Single animal class across four splits\\
Incomplete positive frames excluded & 795 & Includes 85 original test frames; not relabelled empty\\
Boxes clipped/relabelled & 0 & No coordinate corrections recorded\\
Background images retained & 476 & Empty labels in every split\\
Visual verification & 100 images & Saud; all recorded verified. Sample review is not exhaustive annotation proof\\\bottomrule
\end{tabularx}\par
\begin{center}\begin{tabular}{lrrrr}\toprule
Split & Images & Animal boxes & Empty & Excluded frames\\\midrule
Train & 3,013 & 2,819 & 319 & 521\\
Selection-validation & 1,208 & 1,124 & 91 & 129\\
Calibration & 467 & 483 & 13 & 60\\
Test & 554 & 508 & 53 & 85\\\bottomrule
\end{tabular}\end{center}
The calibration partition remains separate. ``Complete'' refers to the annotation ledger, not perfect labels. Metrics on 554 retained test frames do not characterize all 639 original test frames.
\fig{box_verification.jpg}{1}{Human boxes rendered after normalized-coordinate conversion on training examples. Stored images up to 1024 pixels; reference-only illustration, confidence/NMS not applicable. Training input: 640. Source: project CCT20 subset.}

\subsection{Fine-tuning}
YOLO11n starts from \texttt{yolo11n.pt}. The supplied checkpoint embeds a 40-epoch log: input 640, batch 8, seed 42, patience 10, optimizer \texttt{auto}, AMP enabled and box/cls/DFL loss gains 7.5/.5/1.5. Augmentations include horizontal flip .5, HSV jitter, translation .1, scale .5 and mosaic 1.0 closed for the last 10 epochs. These provide appearance/scale variation without proving new-site robustness.
\par\noindent\begin{tabularx}{\linewidth}{@{}l Y r r@{}}\toprule
Run & Selection & Val. AP50 & Val. AP50--95\\\midrule
YOLO11n / 40 epochs & Highest native validation AP50--95, epoch 36 & .90855 & .52302\\\bottomrule
\end{tabularx}\par
Only one trained detector is evidenced; no second-family comparison is claimed. The curve is bound to the served checkpoint hash. Its original training arguments specify GPU device 0; the exact GPU model and original training-run W\&B URL are not archived. The linked W\&B run is an evidence backup.
\fig{detection_training.png}{1}{Train/selection-validation detector curves, 640 pixels. Native training-validator settings, not deployed confidence .045/NMS .30. Epoch 36 maximizes recorded AP50--95. Later training-loss decreases do not improve that selection metric.}

\subsection{Threshold sweep and operating point}
The coarse grid tested confidence $\{.01,.05,.10,.25,.50\}$ and NMS IoU $\{.3,.5,.7\}$. Validation refinement tested .015--.045 at NMS .3. The rule maximizes empty removal subject to frame animal recall $\geq .99$. The selected .045/.30 achieves validation frame recall .990152 and empty removal .285714; .05 fails the recall constraint (.989257). The refined row was supplied during experimentation; the archived sweep CSV contains the coarse grid.
\begin{center}\begin{tabular}{rrrrrr}\toprule
Conf. & NMS & Box P & Box R & Box F1 & Animal-frame recall\\\midrule
.01 & .30 & .4805 & .9306 & .6337 & .9973\\
.05 & .30 & .6634 & .9155 & .7693 & .9893\\
.10 & .30 & .7436 & .9030 & .8156 & .9758\\
.25 & .30 & .8547 & .8790 & .8667 & .9320\\
\textbf{.045} & \textbf{.30} & \textbf{.6513} & \textbf{.9155} & \textbf{.7611} & \textbf{.9902}\\\bottomrule
\end{tabular}\end{center}
\fig{threshold_sweep.png}{.85}{Coarse sweep on 1,208 retained selection-validation frames, input 640 and NMS .30. The dotted line marks the separately refined confidence .045. Box recall and animal-frame recall have different denominators.}

\subsection{Test results}
Native AP uses confidence floor .001, separately from deployment. Operating precision/recall/F1 use .045/.30 and the notebook's one-to-one box matching at IoU $\geq .5$.
\begin{center}\begin{tabular}{lrrrrrr}\toprule
Class & GT & AP50 & AP50--95 & P@op. & R@op. & F1@op.\\\midrule
Animal / all classes & 508 & .89574 & .56954 & .72581 & .88583 & .79787\\\bottomrule
\end{tabular}\end{center}
There are 450 matched TP, 170 FP and 58 FN. Of 501 animal frames, 489 have a detection (97.60\%); 28/53 empty frames are removed (52.83\%). Test animal recall fails .99 despite satisfying the validation constraint.
\begin{center}\begin{tabular}{lrrrr}\toprule
Appearance group & Frames & Box P & Box R & Box F1\\\midrule
Grayscale proxy & 269 & .8367 & .9389 & .8849\\
Colour proxy & 285 & .6258 & .8293 & .7133\\\bottomrule
\end{tabular}\end{center}
Grayscale appearance is not verified nighttime/infrared ground truth. The archived size proxy is normalized box area times $640^2<32^2$, not exact COCO area after letterboxing. All five such small objects are missed. Medium/large are merged because separate results were not exported.
\begin{center}\begin{tabular}{lrr}\toprule
Size proxy & Instances & Recall@op.\\\midrule
Small & 5 & .0000\\
Medium + large & 503 & .8946\\\bottomrule
\end{tabular}\end{center}
The combined recall is derived as $450/503$. Five small instances are insufficient for a stable population estimate.

\subsection{Error analysis}
Saleh reviewed 12 test-gallery frames with 11 FP and two FN counts; one localization error contributes both. This breakdown is for the gallery only, not all 170 FP. Wrong-species classification is not applicable to the single-class detector.
\begin{center}\begin{tabular}{lrr}\toprule
Reviewed FP pattern & Count & Share of 11 gallery FP\\\midrule
Background / foreground artifacts (some uncertain) & 8 & 72.7\%\\
Duplicate / body-part prediction & 1 & 9.1\%\\
Instance/annotation ambiguity & 1 & 9.1\%\\
Localization mismatch & 1 & 9.1\%\\\bottomrule
\end{tabular}\end{center}
A full-test FP taxonomy was not exported. Reviewed causes include rocks, bright branches, glare, blur, small animals and tail-extent disagreement. Possible unannotated animals must be checked before labelling hard negatives.
\fig{detector_failures.jpg}{1}{Reviewed test failures, input 640, confidence .045/NMS .30. Green references and red predictions. This qualitative gallery is not a random frequency sample.}
\textbf{Failure case.} Severe blur in \texttt{58678cd3} accompanies an oversized false-positive box over most of the scene, subsequently harming segmentation. Future work is reviewed hard negatives and quality-aware prompt support, chosen on validation.

\subsection{Inference on images and video}
The implemented interface processes image folders/batches and event metadata, matching Project 5. The 6.70 images/s detector-forward rate is not video FPS or application throughput. No detector-on-video execution or streaming-FPS evidence is supplied. A separately recorded HTTP demonstration remains a presentation deliverable; JSON request evidence is not a video recording.

\section{Segmentation (M5)}
\textbf{Task.} Animal instance masks support per-image union foreground. Relative visible-body area is $|M|/(WH)$; it is not physical size, identity or calibrated distance. No distance estimate is claimed.

\subsection{Ground-truth masks}
\par\noindent\begin{tabularx}{\linewidth}{@{}l Y@{}}\toprule
Item & Evidence\\\midrule
Source & SAM ViT-B proposals from source boxes, followed by recorded team review\\
Images & 150 train / 40 selection-validation / 40 test; 230 total\\
Class & Animal foreground; stored background is zero\\
Review & Yousef: all 230 recorded verified; 34 flagged in the earlier pass before final verification\\
Format & Normalized YOLO polygons; each line is an instance, rasterized at stored image dimensions\\
Training-only cleanup & Six suspect background polygons removed from four images in a derived training copy; original/evaluation labels preserved\\
Remaining concern & Some outlines self-touch; pixel-mask evaluation preserves the frozen release. Visual review identifies suspected reference-only background regions\\\bottomrule
\end{tabularx}\par
Review completion is a ledger status, not a guarantee of perfect references. SAM-assisted labels may favour SAM-like contours. No test labels were corrected to raise reported scores.
\fig{mask_examples.jpg}{1}{Reviewed reference/prediction comparisons on the 40-image development test. Stored resolution up to 1024, inference 640. Green references, red predictions, blended overlap. Primary mask confidence .40/SAM quality .75; v2 fallback below. Reference-only fragments are disclosed.}

\subsection{Model and training}
YOLO11n-seg is continued from the original last checkpoint for 12 epochs on the derived 150-image training set: 640 pixels, batch 4, nominal batch 8, AdamW .0002, cosine decay, no warm-up or mosaic, horizontal flip .5, translation .05, scale .15 and reduced HSV jitter. CPU bfloat16 autocast and FP32 optimizer updates are used. BatchNorm running statistics updated; an attempted callback did not keep them frozen.
\par\noindent\begin{tabularx}{\linewidth}{@{}Y l r@{}}\toprule
Run / change & Split & Mean union IoU\\\midrule
Original M5 last, confidence .25 & Same 40 validation images & .48770\\
Continued M5 alone, confidence .20 & Same validation & .63251\\
Continued M5 + SAM refinement & Same validation & .68396\\
M4 boxes + M5 interior points + SAM & Same validation & .80683\\
Recovery v2, single-mask empty-output fallback & Same validation & .82397\\\bottomrule
\end{tabularx}\par
The baseline is an earlier fine-tuned model, not untouched pretrained weights. Additional training and cleanup are jointly changed; their individual causal effects are not isolated. Epoch 9 maximizes continuation native mask AP50--95. Standalone validation AP50/AP50--95 is .76640/.48899 at low confidence; these are not AP of the three-stage pipeline.
\fig{segmentation_training.png}{1}{Continuation train/selection-validation curves, 150/40 images, 640 pixels and 12 continuation epochs. Native mask-validator AP selects epoch 9. The standalone overlap comparison uses confidence .20, distinct from low-floor AP evaluation.}
\fig{segmentation_ablation.png}{.95}{Incremental comparison on the same 40 frozen validation references, input 640. Rows use their documented policies; this is not a factorial ablation. Final v2 uses M4 primary .40/SAM .75 with fallback floor .045. Repeated selection on 40 images can be optimistic.}

\subsection{Thresholds and post-processing}
M4 proposes at .045/NMS .30, maximum 10 prompts. M5 predicts at confidence .05/NMS .70. A predicted M5 mask is matched to an M4 box at box IoU $\geq .30$; the maximum interior-distance pixel clipped to that box supplies a positive SAM point. The box centre is used if no usable mask exists. Only predicted boxes/masks are inference prompts, not reference annotations.
\par\noindent\begin{tabularx}{\linewidth}{@{}l Y l@{}}\toprule
Setting & Validation comparison & Final\\\midrule
Primary M4 mask confidence & .045, .05, .10, .15, .20, .25, .30, .40, .50 & .40\\
SAM predicted quality & 0, .50, .75 & .75\\
SAM input & 640 vs 768; M5-guided 768 IoU .6954 is below M4-guided 640 & 640\\
Empty-output fallback & Best/all prompt recovery tie at validation IoU .82397; one prompt limits extra outputs & Best eligible prompt\\\bottomrule
\end{tabularx}\par
Primary masks need M4 confidence $\geq .40$ and SAM quality $\geq .75$. When none pass, retain the single highest-confidence mask satisfying the pre-existing M4 floor .045 and unchanged SAM gate. Low-confidence output can still be false positive. SAM quality is a predicted score, not observed IoU or calibrated probability; scores can exceed 1.

Recovery was proposed after v1 test diagnosis: ten images had SAM outputs rejected by the .40 detector gate. Candidate strategies were compared on validation, but the test had already informed development. Accordingly, v2 test performance is not an untouched holdout claim.

\subsection{Results and application measure}
For each image, all animal masks are unioned:
\[\mathrm{IoU}_i=\frac{|P_i\cap G_i|}{|P_i\cup G_i|},\qquad \mathrm{Dice}_i=\frac{2|P_i\cap G_i|}{|P_i|+|G_i|}.\]
Scores are arithmetic image means, not pixel-weighted dataset IoU or a mean over species. All 40 test images, including failed masks, remain in the denominator.
\begin{center}\small\begin{tabular}{llrrr}\toprule
Model & Split & Mask AP50--95 & Mean IoU & Mean Dice\\\midrule
Original M5 last & Validation & .22308 & .48770 & .57919\\
Continued M5 alone & Validation & .48899 & .63251 & .70337\\
Recovery v2 & Validation & Not measured & .82397 & .88895\\
M4-guided v1 & Development test & Not measured & .58022 & .62024\\
Recovery v2 & Development test & .59847* & .73988 & .80008\\\bottomrule
\end{tabular}\end{center}
\textbf{*Operating-policy AP.} COCO instance-mask AP of accepted deployed outputs gives AP50 .80781 and AP50--95 .59847, on 45 reference instances/40 predicted instances in 40 images. Ranking uses M4 confidence. Gates and image-level fallback truncate the outputs: this is not a low-confidence native YOLO/SAM AP sweep, and is not interchangeable with standalone validation AP. Reference hashes were verified.

Mean absolute relative-area error is .023612 of image area (2.361 percentage points), without physical calibration. Nine images improve IoU and none worsen relative to v1; 30/40 reach .70. Three have zero overlap, including two without detector proposals. Ten images use fallback; one still has zero overlap.

\subsection{Error analysis}
Mohannad confirmed AI-assisted visual analysis of 12 lower-IoU cases. Two exceed .70 and are partial successes, not complete failures. The following groups overlap and are not additive.
\par\noindent\begin{tabularx}{\linewidth}{@{}Y l Y@{}}\toprule
Pattern & Examples & Observation\\\midrule
No detector proposal & 2 & Shadowed cat; image-edge squirrel\\
Wrong fallback target & 1 & Tiny annotated squirrel missed; foreground rock/log-like region chosen\\
Oversized mask & 1 & Blur in \texttt{58678cd3}; predicted union occupies 89.23\%\\
Reference-quality concern & 4 & Unclear boundaries or apparent reference-only rock/background regions\\
Partial low-light contours & Several & Body/leg boundaries and edge truncation disagree\\\bottomrule
\end{tabularx}\par
\fig{segmentation_failures.jpg}{1}{Development-test failures at 640-pixel model input; stored images up to 1024. M4 proposal .045/NMS .30, primary confidence .40/SAM .75 with v2 fallback. Green reference/red prediction. Every example remains in the metrics.}
\textbf{Failure case.} In \texttt{58678cd3}, an extra mask covers nearly the entire blurred frame despite a high predicted SAM quality. In \texttt{588c129e}, an apparent reference-only component lies on a rock outside the annotated cat box. This is a reference concern, not a reason to train on background. Any future adjudication must be versioned; present labels/scores are unchanged.

\section{System integration and deployment (M8)}
\subsection{End-to-end pipeline}
CPU PyTorch is served for all learned stages. Counts use every M4 output at .045/.30; mask acceptance does not remove counted boxes. Species preprocessing matches Phase 1: 5\% crop margin, BGR-to-RGB, square resize and ImageNet normalization, with recovered temperature $T=1.03510$. Species confidence below .60, low detector confidence or missing masks triggers review. Mask boxes are associated one-to-one to detector boxes at IoU $\geq .50$.
\begin{center}\begin{tabular}{lrr}\toprule
Measured stage & p50 ms & p95 ms\\\midrule
Detection + image decode & 183.806 & 261.105\\
Crop preprocessing + classification & .003 & 293.215\\
Segmentation + association & 243.090 & 8680.733\\\midrule
Direct full pipeline & 482.640 & 9071.338\\\bottomrule
\end{tabular}\end{center}
These are composite stages, not isolated preprocessing/postprocessing microbenchmarks. Total is measured directly, not a sum of percentiles. Classification p50 is nearly zero because three of the five cases have no detections, skipping crop inference. Segmentation includes repeated M4 execution.
\fig{pipeline_latency.png}{1}{Colab CPU, batch 1; three warmups and 20 sequential replays of five fixed validation cases. Inputs: 640 detector/masks and 224 crops; unchanged thresholds. Positive cases take about 7--9 s versus .35--.53 s for empty/missed cases. Image read included, HTTP excluded. This functional-case mixture is not representative production traffic.}

\subsection{Release and manifest}
\par\noindent\begin{tabularx}{\linewidth}{@{}l Y@{}}\toprule
Item & Served configuration\\\midrule
Detector & YOLO11n; SHA256 prefix \texttt{ac38aaffe735}; one \texttt{animal} class\\
Segmenter & Continued YOLO11n-seg; prefix \texttt{f19fbfb60b2b}\\
Species & New recovery checkpoint; prefix \texttt{b68c7fbdda6b}; original ordered 14-class vocabulary\\
SAM & Official \texttt{sam\_b.pt}; prefix \texttt{ec2df6273261}; verified public weight, not retrained\\
Policy & v2 prefix \texttt{1ec858990417}; paths, full hashes, input sizes, thresholds, fallback and reference identities\\
Versions & Deployment log: Python 3.13.16, PyTorch 2.14.1+cu130 used on CPU, Ultralytics 8.3.221, ONNX 1.23.2. Classifier workflow pins Transformers 4.49.0. Exact ORT version is not exported in this bundle\\
Backup & Curated Drive ZIP includes the three learned checkpoints and policies. Public pretrained SAM is referenced by URL/hash\\\bottomrule
\end{tabularx}\par
Release settings are read from files rather than reselected during service calls; model/helper hashes are checked. Full hashes and metric-source records accompany the report. W\&B v0 backs up an earlier evidence snapshot, not proof that later report/video files were already uploaded.

\subsection{Input/output contracts and HTTP service}
\par\noindent\begin{tabularx}{\linewidth}{@{}l Y@{}}\toprule
Contract & Definition\\\midrule
Request & Multipart \texttt{POST /batch}: repeated \texttt{images} plus JSON \texttt{metadata}, one unique image ID/event record per file; timestamp optional\\
Image processing & Decode BGR; YOLO 640 letterbox; RGB 224 normalized species crops; SAM receives image and predicted prompts\\
Frame output & IDs/time, normalized $xyxy$ boxes, species/score, detector score, review flags, matched mask polygons, relative area and per-stage seconds\\
Event output & Earliest retained timestamp; maximum simultaneous count; per-species maxima. These species maxima need not sum to the event total. No identity tracking\\
Invalid input & Malformed metadata, count/ID errors and empty/unreadable uploads yield HTTP 400. Missing form fields use framework validation. No detections is a valid zero-count response\\
Hosting & FastAPI/uvicorn on a dynamically bound loopback port; instance-specific \texttt{GET /health}. Local prototype, not public hosting\\\bottomrule
\end{tabularx}\par
Example request, substituting the actual port:
\begin{quote}\small\ttfamily
curl -X POST -F 'images=@frame.jpg' -F 'metadata=[\{"id":"frame","event":"e1"\}]' http://127.0.0.1:PORT/batch
\end{quote}
The first three-image attempt returned HTTP 200 but no detections; two positive-label frames had incomplete box annotations. It did not exercise classification or masks and is preserved as a failure. The revised five-case demo retains those misses and the empty frame, adding two previously detected, annotated validation examples. HTTP 200 is recorded; two objects are classified and matched to two masks, exercising all stages. Request time is 19.33 s including transport. This is functional branch coverage, not a representative accuracy sample or recorded video.
\fig{http_outputs.jpg}{1}{Two positive branches of the executed five-case validation HTTP request. Stored images up to 1024, detector/SAM 640, species 224. Operating thresholds unchanged. Boxes, predicted species and green mask contours come from the actual service. Earlier misses remain in the demo evidence.}

\subsection{ONNX export and parity}
Only the detector is exported: FP32, static $1\times3\times640\times640$, opset 17, no graph NMS or simplification. NMS/application logic stay in Python. Three probes (zero, .5 constant, seeded random) pass $rtol=atol=10^{-3}$; maximum absolute error .0003967285.
\par\noindent\begin{tabularx}{\linewidth}{@{}Y l Y@{}}\toprule
Check & Result & Scope\\\midrule
Raw detector tensors & Pass & Shape $1\times5\times8400$, finite, all three probes within tolerance\\
Post-NMS decisions & Not measured & No matching percentage exported\\
Evaluation-set metric parity & Not measured & Native/ONNX AP equality not claimed\\
FP16 / INT8 parity & Not used & CPU FP32 selected\\
Full-pipeline ONNX parity & Not applicable & Classifier, M5 and SAM remain PyTorch\\\bottomrule
\end{tabularx}\par

\subsection{Benchmark}
Hardware is Colab CPU, reported Intel Xeon at 2.00 GHz, four inference threads. A CUDA-enabled package build does not imply GPU inference. Host hardware can vary between sessions.
\begin{center}\small\begin{tabular}{lrrrrr}\toprule
Scope / backend & p50 ms & p95 ms & p99 ms & Images/s & Peak RSS MB\\\midrule
Detector / PyTorch & 135.99 & 192.22 & 205.69 & 6.696 & 907.10\\
Detector / ONNX & 225.05 & 441.98 & 455.22 & 3.817 & 907.10\\
Full pipeline / PyTorch & 482.64 & 9071.34 & 9086.20 & .288 & 2354.72\\\bottomrule
\end{tabular}\end{center}
Detector rows use 10 warmups and 100 random-tensor forward calls, batch 1/640, excluding decode/NMS/application stages. The full path uses three warmups and 20 fixed validation replays. Tail percentiles from 20 calls are descriptive, not production SLO certification. RSS is cumulative process peak, not per-model incremental memory; equal detector-row RSS does not demonstrate equal model memory.
\fig{backend_benchmark.png}{.75}{Raw-forward detector comparison on CPU FP32, batch 1/640, 10 warmups/100 calls; confidence and NMS are outside measurement. ONNX is slower here, motivating native PyTorch serving.}

\subsection{Verdict against acceptance criteria}
\begin{center}\small\begin{tabular}{lrrl}\toprule
Criterion & Target & Measured & Verdict\\\midrule
Animal-frame recall & $\geq .99$ & .97605 & Fail\\
Empty removal & $\geq .20$ & .52830 & Pass\\
Mask union IoU & $\geq .70$ & .73988 & Pass\\
Event-count MAE & $\leq 1$ & .31944 & Pass\\
Full-pipeline p95 seconds & $\leq 2$ & 9.07134 & Fail\\
Species macro-F1 & $\geq .65$ & .48252 & Fail\\
Species ECE & $\leq .10$ & .05128 & Pass\\\midrule
\multicolumn{3}{l}{\textbf{Unattended deployment decision}} & \textbf{NO-GO}\\\bottomrule
\end{tabular}\end{center}
Counts are evaluated on 288 events with complete supervision on every retained test frame; 47 events are excluded. Expected/predicted event count is the maximum reference-box/detector count. Empty-event errors are included. This is not a unique-individual or complete-original-sequence estimate.
\par\noindent\begin{tabularx}{\linewidth}{@{}l Y Y@{}}\toprule
Failed criterion & Proposed remedy & Cost / verification\\\midrule
Animal recall & More independently grouped small/occluded examples; validate higher resolution or an alternative policy & Annotation/compute; evaluate on new unseen data\\
Species macro-F1 & Improve rare-class and camera coverage; targeted adaptation & New labels; preserve class order/calibration and report per-class results\\
CPU latency & Remove repeated M4 only after parity checks; faster refinement or GPU serving & Engineering/compute trade-off; remeasure accuracy and full latency\\\bottomrule
\end{tabularx}\par
The prototype supports research review with human oversight. Out-of-scope uses include identifying people, precise wildlife-location publication, unattended population estimation and physical body size/distance without calibration. Performance NO-GO is separate from course completion; targets were not weakened to obtain GO.

\subsection{Recovered species classifier}
The new model uses 2,819 train, 1,124 selection-validation and 483 calibration crops. Its backbone is frozen in eval mode; the head uses Adam .001, batch 32, up to 12 epochs/patience 4, original flip/colour-jitter preprocessing, and no class weighting, label smoothing or weight decay. Epoch 12 is selected at validation macro-F1 .36854. New library/hardware versions, zero workers for resumption and removal of the original 900-second cutoff distinguish this execution. Identical reproduction is not claimed.
\fig{species_training.png}{1}{Recovery classifier train/selection-validation curves: 224-pixel crops, 14 classes, seed 42; argmax metrics. Temperature is subsequently fitted on the separate calibration split, not these curves or test.}
On 508 historical test crops: accuracy .742126, macro-F1 .482517, top-3 accuracy .899606 and ECE .051280. Fourteen training classes are averaged even though rodent has zero test support. Several rare classes have zero F1. On calibration data, ECE rises from .06081 to .06899 despite NLL-based temperature fitting; temperature scaling is not claimed to improve every ECE estimate. Fresh-model reload verifies saved-logit parity on two calibration inputs.

\section{Obstacles and responses}
\par\noindent\begin{tabularx}{\linewidth}{@{}Y l Y Y@{}}\toprule
Obstacle & Module & Impact & Response / limitation\\\midrule
Unavailable original classifier weight & M8 & Historical model could not be served & New recovery version; original report/metrics preserved\\
Colab GPU quotas and session loss & M5/M8 & CPU continuation and long classifier training & Drive checkpoints and resumable epochs; honest CPU timing\\
Weak initial segmentation confidence & M5 & Early checkpoint gave no useful masks & Validate last checkpoint, continue training and compare policies\\
Background prompts / imperfect masks & M5 & Oversized masks and ambiguous overlap & Training-only cleanup, versioned review; no test-label edits\\
Exposed test during iteration & M5 & Independent holdout status lost & Label v2 results development; require a separate unseen set\\
ONNX slower on CPU & M8 & No export speed advantage & Retain verified export, serve PyTorch\\
Drive mount conflicts / upload waits & M8 & Backup/receipt delays & Preserve local files, safely mount Drive, verify committed artifact\\\bottomrule
\end{tabularx}\par

\section{Submission evidence and limitations}
The supplied audit passes 35/38 checks. Learned weights, policies, reference/failure reviews, metrics, HTTP execution, timing and W\&B receipt are present. The remaining artifacts at that snapshot are report PDF, final presentation PDF and recorded demo video. This report supplies the first; slides/video must be added and the audit rerun. File presence is not content approval or grading completeness.
\par\noindent\begin{tabularx}{\linewidth}{@{}Y Y@{}}\toprule
Evidence & Status at report freeze\\\midrule
Detection conversion/training/AP/operating metrics & Present; full-test FP taxonomy and separate medium/large recall not exported\\
Small-object / night analysis & Small proxy and grayscale groups present; verified IR/night labels unavailable\\
Masks/training/overlap/applicable AP & Present; reference concerns and development-test reuse disclosed\\
Service/stage timing/release/model card & Present; local functional service, not production hosting\\
ONNX parity & Raw tensor checks present; post-NMS/metric parity unmeasured\\
Image/video outputs & Image gallery present; video execution/recorded fallback not evidenced at freeze\\
Acceptance decision & 4/7 targets pass; unattended deployment NO-GO\\
Open-world extension & Not applicable, four-member team\\\bottomrule
\end{tabularx}\par

\section{Self-assessment}
\par\noindent\begin{tabularx}{\linewidth}{@{}Y r r Y@{}}\toprule
Component & Maximum & Provisional self-score & Known gap\\\midrule
Detection M4 & 15 & 12 & Recall/small objects; complete taxonomy/video evidence\\
Segmentation M5 & 15 & 12 & Reference noise, development reuse and AP scope\\
Integration M8 & 15 & 12 & Latency/species targets and incomplete ONNX decision parity\\\midrule
Total & 45 & 36 & Instructor determines actual grade\\\bottomrule
\end{tabularx}\par
These provisional self-scores are a report self-assessment, not an instructor award or a claim that all requirements passed.

\section*{References and reproducibility}
\begin{enumerate}[leftmargin=1.5em,itemsep=2pt]
\item AI251P, \emph{Capstone Project Guidelines}, Project 5, pp.\ 18--19; supplied Phase 2 report template, four-member variant. The submission deadline is 22 November 2026.
\item Caltech Camera Traps / CCT20, LILA: \href{https://lila.science/datasets/caltech-camera-traps}{LILA Caltech Camera Traps dataset}. Human boxes, CDLA-Permissive attribution. Figures are project evaluation examples; no wildlife GPS is reported.
\item Ultralytics YOLO11, software 8.3.221: \url{https://docs.ultralytics.com/}. Native detector AP confidence floor .001; deployment .045/.30.
\item Kirillov et al., \emph{Segment Anything}, 2023. Official SAM ViT-B weight via Ultralytics assets; full SHA256 in the policy.
\item Microsoft \texttt{resnet-50}: \url{https://huggingface.co/microsoft/resnet-50}. Preprocessing SHA256 prefix \texttt{87e3366dd794}; recovery checkpoint prefix \texttt{b68c7fbdda6b}.
\item Accompanying machine-readable evidence: metrics, per-image/per-event CSVs, model/policy identities, detector checkpoint history, continuation logs, current reviews and input provenance. Reproduction tools are on branch \texttt{phase-2}. The new classifier does not reconstruct the original Phase 1 weights.
\end{enumerate}
\end{document}
